\section{A release is a state-compatible system configuration}\label{sec:C32:01}
Deploying a new model endpoint does not by itself deploy a new agent system. The release also includes tokenizer and serialization rules, prompts, context transformations, memory schemas, embedding indexes, skills, tools, permissions, validators, and runtime behavior. A change can be compatible at the API level while incompatible at the level of meaning. A field that remains a string can change from an absolute path to a workspace-relative path; an old skill may then direct a valid call to the wrong location.

Define a release manifest
\begin{equation}
 \nu=(v_\theta,v_q,v_\delta,v_g,v_M,v_I,v_h,v_T,v_A,v_V),
\end{equation}
with versions for the model, tokenizer, decoder, context constructor, memory schema, index, harness, tools, authorization, and validators. The decoder reference binds sampling configuration and versioned constrained-decoding rules; actual per-call settings are retained alongside the model input when runtime-selectable. A comparison claiming fixed decoding holds that configuration fixed. Parsing, stopping, and retries remain part of the versioned harness. A session is pinned to a manifest or governed by an explicit migration rule. Compatibility is a predicate $\mathsf C(\nu,S)$ over configuration and retained state, not merely a comparison of version strings. It can require schema validation, representation conversion, regenerated embeddings, restored cache dependencies, reauthorized tools, and revalidation of procedures extracted under an older policy.

A deployment dependency graph has an edge $u\to v$ when the interpretation or safety of component $v$ depends on the contract supplied by $u$. Changing a tokenizer can invalidate token offsets, cache keys, and learned interface adapters. Changing an embedding model can invalidate nearest-neighbor comparability. Changing authorization can invalidate queued actions and retrieved procedural instructions. The affected closure is the set of reachable dependents. A release gate should examine that closure instead of treating each changed file as an independent patch.

This graph need not be acyclic in its full behavioral form. A policy writes memory that later conditions the policy; a verifier changes the training data that alters the executor. For deployment planning, use a versioned snapshot and a declared update order to break the operational cycle. Record which old artifacts can be read during transition and which new artifacts can be written. A system that writes new-schema records before all readers are compatible has introduced a mixed-version state that rollback may not be able to interpret.

Current framework documentation illustrates the need to distinguish persistence mechanisms. LangGraph's October 2026 persistence documentation separates thread-scoped checkpointers from cross-thread stores and warns that in-memory implementations do not survive process restarts~\cite{langgraph2026docs}. Those documented roles do not establish an end-to-end durability guarantee. The application must select a persistent backend, define retention, verify recovery, and reconcile external effects independently.

